% !TeX root = ../main.tex

\chapter{绪论}

\section{研究背景与意义}

浮游生物是海洋生态系统中的基础组成部分，其丰度、群落结构和时空分布与营养盐循环、藻华暴发、食物网结构以及海洋生态安全密切相关。对浮游生物进行持续、准确和高效的监测，不仅关系到海洋生态过程的科学认识，也直接服务于海洋环境评估、生态预警和资源管理等实际需求\cite{naselli2023ecosystem}。随着显微成像设备、流式成像平台和深度学习识别技术的发展，浮游生物监测已经由低通量人工镜检逐步转向高通量图像分析\cite{eerola2024survey,pitois2025rapid,schmid2023edge}。在这一过程中，如何利用深度学习分类模型对海洋浮游生物进行更高效、更自动化的监测，已经逐渐成为海洋浮游生物监测领域的重要问题。

然而，面向真实业务场景的海洋浮游生物监测并不是一个理想化的集中式封闭分类任务。一方面，海洋浮游生物监测往往由不同海域、不同科研平台或不同业务单位分别建设和维护本地数据中心。不同中心在采样环境、成像设备、类别分布和标注进度上存在明显差异，原始图像及其采样元信息还可能关联采样位置、设备布放和任务背景等敏感内容，因此跨中心直接汇聚原始数据并不总是可行。对于这类多中心协作场景，如何能够在不直接上传原始图像数据和单个样本特征的前提下利用各中心数据，已经成为海洋监测领域的关键课题。联邦学习通过“数据保留在本地、模型更新在中心间聚合”的方式，为此类多中心协作场景提供了可行方案\cite{mcmahan2017fedavg,kairouz2021advances,yang2019federated}。

另一方面，海洋浮游生物监测具有显著的开放世界特征。即使分类模型已经学习了训练集中定义的已知类别（in-distribution，ID），但真实场景下仍不可避免地出现未见浮游生物类别、鱼卵、鱼尾、气泡和颗粒杂波等样本。如果系统对这些样本给出高置信度误判，将直接影响监测结果的可靠性。因此，在真实部署条件下，仅仅提高 ID 类别上的分类准确率并不足以保证系统可用性，模型还需要具备对未知类别的拒识能力。多中心海洋浮游生物监测所面临的，不只是分类精度问题，更是在原始图像数据不能汇聚的约束下进行分布外样本检测的问题。

基于上述背景，本文关注的核心问题是：在不共享原始图像的多中心海洋浮游生物监测场景下，如何为联邦化的图像识别任务构建可实现、可复用的 OOD 检测方法。对于这一问题，后处理式 OOD 检测具有较好的适用性。对于已经完成联邦训练的分类模型，后处理方法能够在不重新设计训练框架、不引入额外生成模型的前提下提供 OOD 检测能力。这种即插即用与轻量化的特性，使其能够更容易地部署在现有的多中心海洋浮游生物监测系统中。围绕这一思路，研究联邦场景下的后处理 OOD 检测方法，不仅具有明确的理论意义，也具有较强的实际应用价值。

\section{国内外研究现状}

海洋浮游生物图像识别已经从低通量人工镜检逐步发展到自动分类、高通量分析和现场部署\cite{eerola2024survey,pitois2025rapid,schmid2023edge}。相关研究一方面关注显微成像或流式成像条件下的自动识别流程，另一方面也面向边缘部署、实时监测和海上自适应采样等应用方向展开探索\cite{pitois2025rapid,schmid2023edge}。这表明，海洋浮游生物图像识别正在从实验室条件下的离线分类任务逐步走向真实业务场景。然而，与封闭集图像分类不同，真实海洋环境中的样本分布会受到海域差异、成像条件、季节变化和设备差异等多种因素影响，系统面临的并不只是已知类别之间的分类问题，还需要能够识别出开放世界中的未知物种。已有研究已经从数据分布偏移（dataset shift）\cite{chen2025datasetshift}、开放集识别（open-set recognition）\cite{kareinen2025openset} 和分布外（out-of-distribution, OOD）检测\cite{han2025benchmarking} 等角度讨论了海洋浮游生物监测系统在真实部署中的失效风险。这说明，海洋浮游生物识别在实际应用中不仅是分类问题，也是开放世界的识别问题。

在多中心海洋监测场景下，这一问题又进一步受到数据协作方式的约束，而无法采用集中式训练。现有工作普遍将联邦学习视为多中心数据协作的重要框架。FedAvg\cite{mcmahan2017fedavg} 建立了最经典的参数平均范式，其基本思想是各客户端保留原始数据，仅上传模型参数在服务器进行聚合。此后，大量研究围绕其通信效率、非独立同分布数据（non-IID）和隐私保护等问题展开，逐步扩展了联邦学习的理论与应用边界\cite{kairouz2021advances,yang2019federated}。现有研究大体可以分为两类：一类主要关注标准联邦分类训练的准确率、收敛性和系统效率；另一类则更强调异构场景下的隐私保护、个性化适配和应用部署。随着 OOD 问题在真实部署中的重要性不断提高，已有工作也开始关注联邦场景下的 OOD 检测\cite{liao2024foogd}。但现有联邦 OOD 方法往往与训练阶段的额外模块或生成式结构高度耦合，带来额外计算开销和实现复杂度。相比之下，对于已经完成联邦训练的分类模型，后处理式（post-hoc）OOD 检测无需重新设计训练框架，且兼具轻量化与“即插即用”的适配性，更适合作为多中心联邦识别系统的扩展。

后处理式 OOD 检测方法以已训练完成的分类模型为基础，不需要重新训练完整系统，部署成本较低，也更容易与现有任务兼容。现有方法大体可以分为两类：一类主要利用输出空间信息，另一类进一步利用特征空间结构。输出空间方法中，MSP\cite{hendrycks2017baseline} 通过最大 softmax 概率刻画模型置信度，是最经典的 OOD 基线之一；Energy\cite{liu2020energy} 则利用 logits 的 log-sum-exp 构造能量分数，在多个视觉任务上表现出较强竞争力。与这类方法相比，特征空间方法进一步利用中间特征与分类结果之间的几何关系。ViM\cite{wang2022vim} 通过主子空间与残差空间分解，结合特征空间结构与 logits 来刻画样本偏离 ID 分布的程度，在多个视觉任务中表现出较强竞争力。在海洋浮游生物 OOD benchmark 上，Han 等\cite{han2025benchmarking} 对 22 种方法进行了统一比较，结果显示 ViM 在该基准上整体表现突出，并在 Far-OOD 场景中优势较为明显。从方法结构上来看，ViM 依赖全局特征均值与协方差，使其能够适配统计量可聚合的联邦学习场景。

尽管 ViM 为后处理式 OOD 检测提供了较有竞争力的技术路线，但在其实际应用中仍存在一个关键问题，即主子空间维度 $k$ 的设定通常依赖经验选择\cite{wang2022vim}。固定维度方案（fixed-k）虽然实现简单，但不同模型的网络结构和特征维度差异明显，经验式 fixed-k 设定往往缺乏统一的统计依据，并可能在不同模型之间带来适配性差异。此外，当 $k$ 取值过大时，部署阶段需要保存的投影矩阵规模和计算 OOD 分数的开销都会增加，不利于后处理模块的轻量化落地。在高维统计场景中，当特征维度与样本量处于相近量级时，样本协方差的特征值容易受到噪声膨胀影响，因此如何从样本协方差或相关矩阵中估计有效维度，已成为一条独立的重要研究线索。ACT（Adjusted Correlation Thresholding）\cite{fan2022act} 属于这一方向的代表方法，它通过相关矩阵谱修正与阈值判别估计因子数量，从而为高维特征空间中的有效维度选择提供数据驱动依据。就 ViM 而言，这类方法为主子空间维度选择提供了可借鉴的统计工具。

综上，已有研究分别从浮游生物开放环境识别、多中心数据协作、后处理式 OOD 检测以及高维统计选维等方面提供了重要基础：真实海洋监测提出了面向 OOD 样本的检测需求，多中心协作使这一需求需要在联邦框架下讨论，ViM 为已完成训练的分类模型提供了 OOD 检测路径，而高维统计中的有效维度选择方法则为 ViM 的主子空间构造提供了统计学支撑。然而，现有研究仍缺少一种能够面向多中心海洋浮游生物监测的分布外检测方法，使其在不共享原始图像的条件下复用联邦分类模型，并以后处理方式同时实现 OOD 检测与主子空间维度自适应选择。

\section{本文主要工作}

针对上述研究缺口，本文围绕联邦后处理 OOD 检测展开研究。在不共享原始图像和单个样本特征的前提下，将 ViM 所需统计量的计算改写为联邦充分统计量聚合过程，构造 FedViM；同时，针对其主子空间维度选择问题，引入 ACT 构造 ACT-FedViM，以实现具有统计学依据的自适应选择。具体而言，本文主要完成了以下工作：

\begin{enumerate}
  \item 面向多中心海洋浮游生物监测场景，将 ViM 引入联邦学习框架，提出了 FedViM。各客户端仅上传一阶与二阶特征充分统计量，服务器据此完成全局均值与协方差重构，在不共享原始图像和单个样本特征的条件下获得 ViM 所需的全局统计量。
  \item 针对 FedViM 主子空间维度的选择问题，提出了 ACT-FedViM 方法，用统计驱动的方式代替固定的维度选择，在保持 OOD 检测性能的同时进一步压缩主子空间维度，增强对不同分类模型的适应性。
  \item 在五个 CNN 模型上对 MSP、Energy、FedViM 和 ACT-FedViM 的 OOD 检测性能进行了系统评估。
\end{enumerate}

全文结构如下：第 2 章给出 FedViM 与 ACT-FedViM 的方法设计；第 3 章介绍数据集、联邦设置、实现细节与评估指标；第 4 章展示五个 CNN 模型上的实验结果，并对 FedViM 与 ACT-FedViM 的性能进行分析；第 5 章给出全文结论，并讨论当前工作的局限与后续研究方向。
